% Einheit 3 von 3 – Pythagoras in Figuren und Körpern (bank/pythagoras/e3.jsonl, zusammenbau v0.1)
%% TODO Zweigzeile Teil 1: was hier gelernt wird (Satz in Schülersprache aus dem Einheitstitel) – Einheit 3
\einheitenkopf[e3]{Einheit 3 von 3 · Pythagoras in Figuren und Körpern}
\zweigzeile{ab Kl. 8, je nach Buch bis Kl. 9 · P10 oft}
\setcounter{aufgabe}{20}

%% TODO Titel: Ich-kann-Satz fehlt in der Bank (Platzhalter: Kettenname) – Nr. 21
%% TODO Anweisung: ein Satz über der ersten Teilaufgabe fehlt in der Bank – Nr. 21
\begin{aufgabe}{Ganz oder halb?}
\begin{teile}
\teil Kegel mit Durchmesser $18$ cm – welche Länge gehört ins Teildreieck aus Höhe, Radius und Mantellinie? \\ \kreuz{$18$ cm} \\ \kreuz{$9$ cm}
\end{teile}
\end{aufgabe}

%% TODO Titel: Ich-kann-Satz fehlt in der Bank (Platzhalter: Kettenname) – Nr. 22
%% TODO Anweisung: ein Satz über der ersten Teilaufgabe fehlt in der Bank – Nr. 22
\begin{aufgabe}{Figuren und Körper}
\swa{Raute mit Diagonalen: ein rechtwinkliges Teildreieck nachfahren. Welche Seiten sind Katheten, welche ist die Hypotenuse?}{\raute[diagonalen]{4}{2.5}}
\swz{Gleichschenkliges Dreieck: Grundseite $18$ cm, Schenkel $41$ cm – Höhe $h$?}{\dreieck{(0,0)}{(2.02,0)}{(1.01,4.5)}{$41$ cm}{}{$18$ cm}{}{}{}}
\swz{Gleichschenkliges Dreieck: Grundseite $22$ m, Schenkel $61$ m – Höhe $h$?}{\dreieck{(0,0)}{(1.65,0)}{(0.82,4.5)}{$61$ m}{}{$22$ m}{}{}{}}
\swz{Gleichschenkliges Dreieck: Grundseite $56$ cm, Schenkel $53$ cm – Höhe $h$?}{\dreieck{(0,0)}{(4.5,0)}{(2.25,3.62)}{$53$ cm}{}{$56$ cm}{}{}{}}
\swz{Gleichschenkliges Dreieck: Grundseite $48$ mm, Schenkel $30$ mm – Höhe $h$?}{\dreieck{(0,0)}{(4.5,0)}{(2.25,1.69)}{$30$ mm}{}{$48$ mm}{}{}{}}
\swz{Gleichschenkliges Dreieck: Grundseite $60$ cm, Schenkel $50$ cm – Höhe $h$?}{\dreieck{(0,0)}{(4.5,0)}{(2.25,3)}{$50$ cm}{}{$60$ cm}{}{}{}}
\swfrage{Was bleibt gleich, was ändert sich?}
%% TODO Darstellung für schwach fehlt in der Bank (pythagoras-e3-k2-s2-v1; Abschnitt „Für schwache Schüler“)
\swz{Gleichschenkliges Dreieck: Grundseite $24$ cm, Höhe $35$ cm – Schenkel?}{}
\swz{Gleichschenkliges Trapez: parallele Seiten $50$ cm und $32$ cm, Höhe $40$ cm – Schenkel?}{\trapez[hoehe=h]{5}{3.2}{4}}
%% TODO Darstellung für schwach fehlt in der Bank (pythagoras-e3-k2-s4-v1; Abschnitt „Für schwache Schüler“)
\swz{Parallelogramm: schräge Seite $5{,}3$ m. Der Fußpunkt der Höhe liegt auf der Verlängerung der Grundseite, $2{,}8$ m hinter der Ecke – Höhe?}{}
%% TODO Darstellung für schwach fehlt in der Bank (pythagoras-e3-k2-s5-v1; Abschnitt „Für schwache Schüler“)
\swz{Dreieck ABC mit stumpfem Winkel bei B. Die Höhe von C trifft die Verlängerung von AB im Punkt D. $CD = 3{,}6$ cm, $BC = 3{,}9$ cm – wie lang ist BD?}{}
%% TODO Darstellung für schwach fehlt in der Bank (pythagoras-e3-k2-s6-v1; Abschnitt „Für schwache Schüler“)
\swz{Rechteck $1{,}8$ m mal $8{,}0$ m – Diagonale?}{}
%% TODO Darstellung für schwach fehlt in der Bank (pythagoras-e3-k2-s7-v1; Abschnitt „Für schwache Schüler“)
\swz{$P(1 | 1)$ und $Q(4 | 5)$ – Länge der Strecke PQ?}{}
%% TODO Darstellung für schwach fehlt in der Bank (pythagoras-e3-k2-s8-v1; Abschnitt „Für schwache Schüler“)
\swz{$E(-4 | -3)$ und $F(5 | 9)$ – Länge der Strecke EF?}{}
\swz{Quadratische Pyramide: Grundkante $18$ cm, Höhe $40$ cm – Seitenhöhe $h_s$?}{\pyramide{4}{4}{8.89}{$18$ cm}{}{$40$ cm}}
\swz{Kegel: Durchmesser $66$ cm, Höhe $56$ cm – Mantellinie $s$?}{\kegel{1.77}{3}{}{$56$ cm}{$s$}}
\swz{Kegel: Mantellinie $82$ cm, Radius $18$ cm – Höhe $h$?}{\kegel{0.68}{3}{$18$ cm}{$h$}{$82$ cm}}
%% TODO Darstellung für schwach fehlt in der Bank (pythagoras-e3-k2-s12-v1; Abschnitt „Für schwache Schüler“)
\swz[3]{Beschreibe ohne Zahlen, wie man die Seitenhöhe einer quadratischen Pyramide aus Grundkante und Körperhöhe berechnet.}{}
\swz{Ein Turm besteht aus einem Zylinder ($18$ m hoch) und einem Kegeldach mit Radius $2{,}0$ m und Mantellinie $5{,}2$ m. Wie hoch ist der Turm?}{\zylinder{1}{3}{}{$18$ m}}
\swz{Quader $2$ dm mal $3$ dm mal $6$ dm – Raumdiagonale?}{\quader{1.33}{2}{4}{}{}{}}
\swz[3]{Ein Pavillon hat ein kegelförmiges Zeltdach mit dem Durchmesser $7{,}2$ m. Die Höhe des Dachs ist bekannt, steht hier aber nicht. Beschreibe, wie man die Länge einer Zeltstange vom Dachrand bis zur Spitze berechnet \hfill (P10 2018 OS).}{\kegel{1.8}{1.5}{}{$h$}{$s$}}
\swz[3]{Ein Getreidesilo besteht aus einem Zylinder mit der Höhe $19{,}5$ m und einem Kegeldach mit Radius $3{,}9$ m und Mantellinie $6{,}2$ m. Berechne die Gesamthöhe des Silos auf eine Dezimale \hfill (P10 2026 FOR).}{\zylinder{1}{3}{}{$19{,}5$ m}}
\swz[3]{Ein zylindrisches Glas ist innen $11$ cm hoch und hat den Radius $2{,}8$ cm. Ein Strohhalm steht schräg von der Bodenkante zur gegenüberliegenden oberen Kante und ragt $4$ cm heraus. Berechne die Länge des Strohhalms auf eine Dezimale \hfill (P10 2022 OS).}{\zylinder{0.76}{3}{$2{,}8$ cm}{$11$ cm}}
\swz[3]{Die Punkte R und S sind im Koordinatensystem eingezeichnet. Berechne die Länge der Strecke RS auf eine Dezimale \hfill (P10 2019 OS).}{\begin{ksys}[xmin=-4,xmax=5,ymin=-4,ymax=4,ablesen] \punkt{-2}{3}{R} \punkt{4}{-1}{S} \end{ksys}}
\end{aufgabe}

%% TODO Titel: Ich-kann-Satz fehlt in der Bank (Platzhalter: Kettenname) – Nr. 23
%% TODO Anweisung: ein Satz über der ersten Teilaufgabe fehlt in der Bank – Nr. 23
\begin{aufgabe}{Rampe, Leiter, Seil an der Wand mit Skizze}
\swz{Ein Seil ist an einer Wand in $5{,}6$ m Höhe befestigt und $3{,}3$ m vor der Wand im Boden verankert. Wie lang ist das Seil?}{\dreieckrw{2.95}{5}{$3{,}3$ m}{$5{,}6$ m}{?}}
\end{aufgabe}

%% TODO Titel: Ich-kann-Satz fehlt in der Bank (Platzhalter: Kettenname) – Nr. 24
%% TODO Anweisung: ein Satz über der ersten Teilaufgabe fehlt in der Bank – Nr. 24
\begin{aufgabe}{Figuren und Körper – Fehler finden, Begründen, Anwendung, Darstellungswechsel}
\swz[3]{Gleichschenkliges Dreieck: Grundseite $1{,}4$ m, Schenkel $2{,}5$ m. Nils rechnet: \rechnung{h^2 &= 2{,}5^2 - 1{,}4^2 \\ h^2 &= 4{,}29 \\ h &\approx 2{,}1 \text{ m}} Wo ist der Fehler?}{}
\swfrage{Warum rechnet man im gleichschenkligen Dreieck mit der halben Grundseite?}
\swz[3]{Ein Regenschirm ist $62$ cm lang und lässt sich nicht kürzer machen. Passt er flach in einen Koffer, der innen $55$ cm lang und $35$ cm breit ist?}{}
\swa{Zeichne $A(-2 | -1)$ und $B(4 | 3)$ ein, dazu das rechtwinklige Dreieck mit waagerechter und senkrechter Kathete. Beschrifte die Kathetenlängen.}{\begin{ksys}[xmin=-3,xmax=5,ymin=-2,ymax=4] \end{ksys}}
\end{aufgabe}

\uebersichtskasten{Teildreieck: Suche das rechtwinklige Dreieck in der Figur, fahre es nach und benenne neu – Hypotenuse gegenüber dem rechten Winkel, die anderen beiden sind Katheten; oft ist eine Kathete nur die halbe Seite, die halbe Diagonale oder der Radius. \\ Gleichschenkliges Dreieck, Grundseite 16 cm, Schenkel 17 cm: h² = 17² − 8² = 289 − 64 = 225 \quad h = 15 cm \\ Koordinaten: Die Katheten sind die Unterschiede der x-Werte und der y-Werte, die Strecke ist die Hypotenuse. \quad A(2 | 1), B(8 | 9): AB² = 6² + 8² = 100 \quad AB = 10 \\ Körper: Höhe, Radius (oder halbe Grundkante) und Mantellinie (oder Seitenhöhe) bilden ein rechtwinkliges Dreieck; die schräge Strecke ist die Hypotenuse. \quad Kegel r = 20 cm, h = 21 cm: s² = 400 + 441 = 841 \quad s = 29 cm}
